\documentclass[../main.tex]{subfiles}
\graphicspath{{\subfix{../images/}}}
\begin{document}

\chapter{Proposal}

This chapter introduces our proposed approach to the lightweight FER problem, as informed by our literature review process. 
We outline the research problem, progressing from general FER methods to more optimized, lightweight solutions. 
A comprehensive description of our proposed architecture and training procedure is provided. 
Additionally, we review related approaches in lightweight FER and FER with self-attention mechanisms. 
Lastly, we justify the necessity for such a research project based on current evidence.
\newpage

\section{Problem definition}

Facial expression recognition (FER) systems provide automated solutions that analyze facial features to identify the emotion being conveyed. 
In the context of static or image-based FER, feature-extracting CNNs followed by fully-connected  classifiers have been the generally predominant approach \cite{canal2022survey}.
Following the wider trend in deep computer vision, networks with self-attention mechanisms have been applied to the FER task to obtain significant classification improvements \cite{lengwe2023pattlite, apvit2023, xue2021transfer, ma2023asf, huang2021, pecoraro2022lhc}.
In particular, Vision Transformers are quite adept at the FER task when training with ITW conditions.
This is, in part, thanks to the network's ability to pay greater attention to more important facial regions, which mitigates the effect of occlusion and pose variations \cite{karnati2023survey}.\\

While several of these approaches have achieved state-of-the-art results, many of them incur a heavy resource cost in terms of network parameters and floating point operations per-second (FLOPs).
When considering that the feature extraction step is the system's bottleneck \cite{karnati2023survey}, these solutions can become unsuitable for resource-constrained environments.
Some FER researchers have developed lightweight approaches that aim to maintain a comparable level of performance while greatly reducing parameters and computational cost \cite{eff-face2021, lengwe2023pattlite, hewitt2018, mitre2020, barros2020fc}.
These proposals generally employ well-studied, lightweight CNN architectures and combine them with custom modules or training strategies to maintain robustness for FER ITW.
Such methods have achieved great results, with some even setting the current state-of-the-art \cite{lengwe2023pattlite}.
However, this is still only a small subset of FER research, and only an even smaller subset has been tested on ITW databases \cite{eff-face2021}.\\

One such lightweight FER network is EfficientFace \cite{eff-face2021}.
This network achieves great accuracy results on the RAF-DB (88.36\%) and AffectNet (63.7\%) datasets when taking into account its very low parameter count of 1.28 million.
However, other Transformer models like APViT \cite{apvit2023} are capable of achieving much better results (RAF-DB 91.98\%, AffectNet 66.91\%), albeit with a significant increase in performance cost (65.2 M parameters).

% Typically studied under Ekman and Friesen's \cite{ekman1971} basic emotions model, expressions can be mapped to 6 basic emotions: anger, disgust, fear, happiness, sadness, and surprise.
% Current research also adds neutral and sometimes a contempt category \cite{karnati2023survey}.
% Under this paradigm, automated FER can be taken as a multi-class image classification task.\\

% Early researchers applied classical image processing approaches such as hand-crafted features coupled with support vector machines or dynamic Bayesian networks for suitable classification performance \cite{canal2022survey}.
% However, with the advancements in neural network research, processing power, and data gathering capabilities, deep learning-based approaches have become the norm in FER \cite{karnati2023survey, canal2022survey, li2022survey}.\\

% The other important change in FER research over the years is the shift from lab-controlled data to in-the-wild (ITW) datasets.
% While early FER datasets were small in both sample and subject size, newer ITW sets provide thousands of images to train and evaluate with \cite{li2022survey}.
% The current biggest dataset, AffectNet \cite{affnet2019}, contains more than 1 million facial images.
% Although this orders of magnitude increase in the available data is great for robust training, ITW datasets are far more challenging, even for current deep learning methods.
% Whereas head pose, occlusion by hair or accessories, and illumination were strictly controlled in lab datasets, ITW databases impose no such restrictions.
% This leads to vastly degraded performance of FER systems when tested with ITW conditions \cite{zhao2021manet, ma2023asf}.\\

\section{Proposal} \label{c3:proposal}

In order to address the previously stated problem, we intend to improve EfficientFace's (EF) \cite{eff-face2021} FER accuracy by enhancing its label distribution learning (LDL) component.
To accomplish this, we will modify the label distribution generator (LDG) in three significant ways that leverage different aspects of both EF and APViT \cite{apvit2023}.\\

First, we will replace the LDG's ResNet-50 backbone with IR-50 \cite{deng2019ir50} and load APViT's backbone weights onto the first three stages.
This will enable us to transfer knowledge from a network that has already been pretrained on the FER task.
Second, we will adapt EF's local-feature extractor and channel-spatial modulator modules into the LDG.
With this, the network should have better awareness of the global and local facial features related to expression.
Third, we will train the modified LDG using Zhao et al.'s \cite{eff-face2021} LDL approach, where APViT will act as the distribution generator.
With the evidence \cite{eff-face2021} in favor of LDL training, we expect that applying this method will improve the network's FER accuracy.
After modifying and training the LDG, the final step will be to re-train EF with the authors' original training approach, but using the modified LDG instead of the base LDG.\\

The rationale for this proposal is given by two main intuitions.
First, we essentially apply \cite{eff-face2021}'s proposal at a larger scale by using the modified LDG and APViT as the trained and generator networks respectively.
Doing this should yield a LDG model that achieves an improved accuracy over its original counterpart.
Furthermore, an improved LDG should also improve EF's accuracy given the effectiveness of the LDL method.
Secondly, we argue that the modified LDG will serve as an intermediary that can transfer knowledge smoothly between two networks of vastly different sizes.
In doing this, we are exploiting APViT's superior FER performance and domain knowledge to improve EF's accuracy while keeping its low resource usage.
It is in this same vein that we refrain from modifying the EF network directly, and instead modify the portion (LDG) that is only used in training and not during inference.\\

\section{Related work}

The following section describes various deep learning-based approaches for static FER.
Given the nature of our proposal, we focus on lightweight models and those that implement self-attention mechanisms such as vision transformers.
%We also briefly mention other approaches encountered during the literary review process.

\subsection{Lightweight FER models}

Earlier lightweight FER studies tended to adapt existing CNN modules or architectures towards the FER problem.
Hewitt and Gunes \cite{hewitt2018} performed a comparative study between FER architectures for mobile environments focusing on small model sizes.
The study proposed and compared three network variants based on popular CNN architectures: AlexNet, VGGNet, and MobileNet. 
Mitre-Hernández and Ferro-Perez \cite{mitre2020} proposed ResMoNet, which stacks different optimized CNN blocks for lightweight FER.
Their approach achieves a smaller size and lower FLOPs than MobileNet, with a lower memory utilization than other comparable networks.
FaceChannel, proposed by Barros et al. \cite{barros2020fc}, is a FER model that predicts emotions based on the dimensional, valence-arousal model instead of the categorical one.
It is a small network based on VGG16 that trims down the number of layers to obtain a significantly smaller number of parameters.\\

More recent studies have combined the current versions of lightweight models such as ShuffleNet and MobileNet with custom modules to achieve outstanding performance.
Zhao et al.'s \cite{eff-face2021} EfficientFace integrates specialized modules for local-global feature extraction into the lightweight ShuffleNet-V2 architecture.
Additionally, they adopt a label distribution learning approach where they train EfficientFace to learn the outputs of a FER-trained ResNet-50 instead of the usual dataset labels.
Wu et al. \cite{wu2020edge} also propose a lightweight FER network based on ShuffleNet-V2 intended for a cloud-edge device model.
Their approach proposes a system that runs the lighter network components on edge devices.
However, as an optional step, the intermediate feature maps are sent to a more powerful cloud device that computes channel-spatial attention maps with custom pooling modules.
PAtt-Lite by Le Ngwe et al. \cite{lengwe2023pattlite} instead uses a truncated MobileNetV1 as its backbone.
This truncated network is followed by a patch extraction block that splits the feature maps into four patches for local feature extraction.
Finally, the patch feature maps are fed to a special module that uses self-attention for classification.

\subsection{Self-attention in FER}

To the best of our knowledge, Xue et al. \cite{xue2021transfer} were the first group of authors to apply Vision Transformers (ViTs) to the FER problem.
They propose the TransFER model, which applies a specialized type of dropout at different stages to help the model learn more diverse representations.
The attention dropping modules are applied at the convolutional and ViT encoder stages to randomly drop feature maps or self-attention heads respectively.
Huang et al. \cite{huang2021} presented a network that uses an attention mechanism at both the low and high levels of feature learning.
At the low level, a grid-wise attention mechanism is used to learn long-range biases between different facial features.
At the high level, the feature maps are transformed into visual tokens and fed to a ViT encoder that learns global representations of expressions.
Ma et al. \cite{ma2023asf} instead propose feeding richer feature maps to the ViT encoder network by diversifying the network inputs.
Specifically, they extract the local binary pattern (LBP) from the input image and extract features from both the RGB and LBP data using two parallel ResNet-18 networks.
These separate maps are fused together by a special module that also extracts local and global features and then feeds them to the ViT encoder.
Iterating on the TransFER model \cite{apvit2023}, Xue et al. propose the APViT network.
They employ two attention pooling mechanisms to drop feature map patches and visual tokens based on their attention scores, discarding the less important regions.
In doing this, they manage to significantly lower the model's FLOPs while improving accuracy by avoiding noisy features.\\

A smaller number of authors have also applied the self-attention mechanism for FER without using a ViT encoder component.
Pecoraro et al. \cite{pecoraro2022lhc} propose a special module that applies multi-headed self-attention to the output feature maps of convolutional layers.
This module is intended to be compatible with existing networks so that it only has to be inserted between the CNN layers without requiring further modification.
The attention classifier in PAtt-Lite \cite{lengwe2023pattlite} also employs dot product self-attention on the last fully-connected layers of the network.
The ablation study shows that this attention classifier obtains significant performance improvements when applied.

%\subsection{Other approaches}

\section{Justification}

FER is a field of research with many potential applications for a variety of sectors.
Current research has studied the effectiveness of FER in areas such as healthcare \cite{muhammad2017health}, autonomous vehicles \cite{zaman2022car}, industrial Internet-of-Things \cite{xi2021iot}, education \cite{savchenko2022edu}, and human-robot interaction \cite{wu2021hri}.
Thus, new developments in FER may bring with them important benefits to such related applications. 
Prior state-of-the-art networks such as RAN \cite{wang2020ran}, published in 2020, have been surpassed by current works like APViT \cite{apvit2023} and POSTER \cite{zheng2023poster} by a good margin, which suggests an accelerated rate of improvement.\\

While deep learning architectures have achieved impressive results not only in FER, but also in general computer vision tasks, the resource cost incurred by these state-of-the-art models poses a significant barrier for real world use.
For example, top performing models such as APViT \cite{apvit2023}, POSTER \cite{zheng2023poster}, and POSTER++ \cite{mao2023posterv2} have a number of parameters in the range from 43.7 million to 71.8 million \cite{lengwe2023pattlite}.
Such requirements limit what can be done with FER systems intended to run in resource-constrained environments such as mobile or IoT edge devices.
As such, there is a real need for further research into lightweight FER.
Unfortunately, the community involved in this endeavor remains small with only a few authors across several years focusing on optimized FER \cite{eff-face2021}.


\end{document}